\section{Results on the Live Benchmark}
\label{sec:leaderboard}

Every score in this section was computed by the competition's scoring service
against ground truth we do not hold. We implemented none of the four metrics.
Quantities derived from those scores --- the decomposition, the fit, the ratios
--- are our own arithmetic and are identified as such.

\subsection{A video-blind constant on the official leaderboard}

Table~\ref{tab:lb} reproduces the final private leaderboard as published, with
the video-blind controls below the rule: submitted after the closing date and
scored by the same service against the same private split. They are not
leaderboard entries.

A constant risk score of $0.51$ scores \textbf{2.35291}. That equals the
eighth-placed team's score at the five decimals the leaderboard displays,
exceeds five of the thirteen teams including our own, and falls $0.14874$ short
of the winner. Per the scope stated in Section~\ref{sec:intro}, the only
inference we draw from the coincidence at eighth place is that the metric
assigns those two submissions the same value; a whole family of curves --- every
curve above $0.5$ at frame~0 that never dips --- receives it, and we infer
nothing about what any team submitted.

We quote the difference before the ratio deliberately. The metric is a weighted
sum with no meaningful origin, so a ratio of two scores can be made to say
almost anything by shifting the scale. The difference of $0.14874$ and the
decomposition below are scale-free; that the constant reaches 94.1\% of the
winning score is a convenience, not evidence.

One further row is worth stating plainly. The organisers' own sample
submission, resubmitted byte for byte, scores $1.04203$, while the benchmark row
displayed on the leaderboard scores $0.58333$ --- so those are not the same
file, and the sample submission is not the origin of the published benchmark
score.

\subsection{Decomposing the score}

Because every video-blind submission is identical across clips, all of them
induce the same ordering over clips and earn the same AP and the same AUC,
which therefore cancel in any difference between two of them
(Section~\ref{sec:probes}). Differencing the probes separates the terms;
Table~\ref{tab:decomp} gives both the probe scores and the components recovered
from them.

Read the ratio precisely. It compares the timing terms at their maximum against
the discrimination terms at chance, not maximum against maximum. The maxima
ratio is not identifiable from outside, because the weights are undisclosed,
but it is bounded: the winner's discrimination terms are worth at least
$0.14874 + 0.35105 = 0.49979$, giving a maxima ratio of at most $4.01$. Both
framings say the same thing, and the measured statement is the safer one --- a
submission that reads nothing collects $2.00186$, and the best discrimination
anyone demonstrated on this benchmark added $0.14874$ to it.

Two further readings follow. A constant at $0.51$ and a constant at $0.99$
score identically, so the metric reads only whether the curve stands above the
threshold and never by how much. And the dip probe scores $1.85347$ where the
per-frame stable-timing coefficient predicts $1.79041$, so the additive split of
the timing total into its two parts is close but not exact; we report the split
with that tolerance attached. The total itself is measured directly between two
flat curves and does not depend on the split.

\subsection{The score is linear in the crossing frame}

Over the region in which the alarm still precedes the accident in every clip,
the score falls by $0.016647$ for each frame of delay --- one point per $60.1$
frames, or $0.4994$ per second at 30\,fps (Fig.~\ref{fig:linear}). A dedicated
one-frame experiment confirms the slope directly: a step at frame 75 scores
$1.10445$ and a step at frame 76 scores $1.08779$, a difference of $0.01666$
against $1/60 = 0.016667$.

The line was then asked to predict curves it had not seen. Three logistic
curves crossing at frames 60, 75 and 90 are predicted to within
$6.4\times10^{-3}$, the closest to $6\times10^{-5}$. A fourth, crossing at frame
105, lies beyond the fitted region where the curve has already begun to bend;
the model is not expected to hold there and we report it rather than omit a
held-out point that missed.

\subsection{The published definition does not reproduce the scorer}
\label{sec:anomaly}

One result does not fit, and we report it rather than smooth it. A linear ramp
crosses the threshold at frame 75, exactly as a step function does, and scores
$0.037$ lower. That gap is $2.24$ frames of slope, not an integer, so it cannot
be a crossing-frame effect. We therefore submitted a diagnostic family holding
the crossing frame at 75 and varying only the shape away from it
(Table~\ref{tab:anomaly}).

The four curves at the head of that table cross at the same frame and never
fall back, and the three step-like ones hold identical values at frames 74 and
75. Under the metric as published they must score identically: both timing
terms are functions of threshold crossings alone, and video-blind curves are
identical across clips, so the discrimination terms cannot separate them
either. They span $0.09093$, which is $5.5$ frames of slope.

We do not have an explanation and state the two candidates rather than choose.
Either the discrimination terms are not invariant to curve shape among
video-blind submissions --- which would weaken the cancellation argument of
Section~\ref{sec:probes}, though it is hard to reconcile with a step at frame
140 and a flat $0.49$ scoring identically despite very different frame
orderings --- or the timing terms read something beyond the first crossing that
the three agreeing curves share and the ramp does not.

Either way one conclusion is available without resolving it, and it is a
finding in its own right: \textbf{the benchmark's published formula does not
reproduce its scorer}. An entrant cannot compute this competition's score from
its stated definition even given the labels. The headline result is unaffected,
because the constant and the floor are both flat curves whose difference is
measured directly; what carries a caveat is the split of the timing total.

\subsection{A property of the withheld ground truth, recovered}

For a step at frame $k$, both timing terms equal $\max(t_{ai} - k, 0)$ on each
contributing clip, so the score above the floor is $c\,\mathbb{E}[\max(t_{ai} -
k, 0)]$ with $c$ the sum of the two timing weights. Differentiating, the slope
at $k$ is $-c\,P(t_{ai} > k)$: the slope measures the survival function of the
accident-onset distribution, scaled. The leaderboard is therefore reporting, one
submission at a time, the shape of an annotation no entrant has seen.

The initial segment estimates $c$ at $0.016667$, and dividing the timing total
by $c$ gives a mean onset of $120.1$ frames --- about four seconds into a
five-second clip. This explains mechanically why the constant does so well
here: because the accident is always late in the window, an alarm at frame~0 is
credited with roughly 120 frames of anticipation on every clip.

Two caveats belong with it. Under the model the magnitude of the slope must be
non-increasing in $k$, because a survival function cannot rise; it is not, by
roughly ninety times what five-decimal rounding permits, and we suspect the
same second-order effect as in Section~\ref{sec:anomaly}. And when $k$ exceeds
$t_{ai}$ the maximum in the published definition is taken over an empty set and
is undefined; we assume the scorer takes it as zero. We deliberately do not
separate the weights from the metric values: every quantity above is a product
of the two, and the positive rate cancels identically out of the onset ratio,
so this measurement says nothing about what fraction of the corpus contains an
accident.

\subsection{The case study: our own entry}
\label{sec:casestudy}

We reached this question by entering the competition, with a training-free
ensemble combining CLIP-based semantic scoring, optical-flow motion energy and
a third channel, fused by grid-searched weights with monotone clamping and
temporal post-processing. It placed ninth of thirteen at $2.00585$ ---
\emph{below the video-blind constant}. Auditing it against its released code
rather than against our description of it found three defects of exactly the
kind this paper is about. The third ``semantic'' channel is not a captioning
model: it thresholds the same frame-difference statistic the motion channel
uses and selects one of three fixed strings, so the ensemble has two
independent inputs rather than three --- which is why that channel posts an
\emph{earlier} mean crossover than the full ensemble. The monotone clamp forces
precisely the condition STTA tests, so any report of STTA compliance measures
that the clamp executed. And a temporal resampling stage evaluated the curve at
index $\alpha t$ while emitting it at $t$, so the score reported for frame $t$
was computed from a frame that had not yet arrived; a six-frame gain once
attributed to it is withdrawn, and the released configuration disables the
stage. No measurement above depends on this system. It is here because a
critique of evaluation practice written by authors who had not audited their
own entry would be worth less.
